% Vietnamese translation for OI Public Library
% Source-File: IOI中国国家候选队论文/2003/高正宇/高正宇.ppt
\documentclass[11pt]{article}
\usepackage{oipl-vn}

\title{Bản trình chiếu: Phân loại ví dụ và kỹ thuật giải}
\author{Gao Zhengyu}
\date{2003}

\begin{document}
\maketitle

\origtitle{Slide companion for problem classification examples}
\sourcepath{IOI China candidate team papers / 2003 / Gao Zhengyu / Gao Zhengyu.ppt}

\section{Phạm vi bản dịch}

Tệp gốc chỉ có bản PowerPoint. Phần chữ khôi phục được gồm danh sách nguồn
bài như IOI 2000, IOI 2002 \textit{Rods}, NOI 2002, CEOI 2001
\textit{Chain}, IPSC 2002 \textit{Chamber of Secrets}, và POI \textit{Ulam}.

\section{Mạch trình bày}

Bản trình chiếu có vẻ dùng nhiều bài thi khác nhau để phân loại kỹ thuật:
đếm trạng thái, dựng mô hình, tìm quy luật, và xử lý trường hợp đặc biệt.
Các con số như \(32768\) và biểu thức \(\min(n-m,0)\) xuất hiện như các mốc
hoặc công thức trong ví dụ.

\section{Cách đọc}

Vì phần chữ mô tả chính không khôi phục được, bản LaTeX này ghi lại danh sách
bài và ý nghĩa didactic của slide: dùng nhiều đề nguồn khác nhau để chỉ ra
cách nhận diện cấu trúc thuật toán trước khi cài đặt.

\section{Ghi chú}

Nội dung chi tiết nên được mở rộng khi có bản chuyển đổi slide sang ảnh hoặc
PDF. Hiện tại, việc tự bổ sung lời giải cụ thể cho từng bài sẽ vượt quá dữ
liệu đọc được từ nguồn.

\end{document}
